\section{Metodologia}
\label{sec:metodologia}

\pendente{Descrever o percurso metodológico completo, incluindo classificação
da pesquisa, procedimentos de levantamento de requisitos, processo de
desenvolvimento e forma de avaliação.}

\subsection{Classificação da pesquisa}

\pendente{Natureza, objetivos e procedimentos técnicos, com a fonte
metodológica que sustenta a classificação.}

\subsection{Levantamento de requisitos}

\pendente{Entrevistas com servidores do setor de licitações e análise de
documentos reais de processos já encerrados. Registrar em
\code{TCC/requisitos/} cada sessão: o que foi dito, o que foi inferido e o que
ficou em aberto.}

\subsection{Tratamento dos dados de domínio}

Documentos reais de processos licitatórios não são versionados nem reproduzidos
no texto. Cada documento recebido é mantido fora do repositório e dele se deriva
uma versão sintética, com a mesma estrutura e nenhum dado verdadeiro, que é a
única utilizada em testes, exemplos e ilustrações deste trabalho.

\pendente{Detalhar o procedimento de derivação e discutir suas limitações para
a avaliação da solução.}

\subsection{Arquitetura da solução}

\pendente{Esta subseção recebeu, por decisão de 24/08/2026, o conteúdo técnico
que antes ocupava o referencial teórico. O motivo: recuperação de informação e
representação vetorial não são base teórica emprestada de terceiros neste
trabalho --- são a contribuição, e o modelo do curso é explícito em não admitir
nova revisão bibliográfica dentro do desenvolvimento.}

\subsubsection{Recuperação de informação sobre a legislação e os modelos}

\pendente{Recuperação aumentada por geração aplicada à legislação e aos modelos
de documento já aprovados. Descrever a estratégia de fragmentação, o modelo de
representação escolhido e o critério de recuperação. Há material reaproveitável
dos fichamentos da Iniciação Científica, mas cada citação precisa ser
reconferida contra o PDF.}

\pendente{Posicionar contra o contra-argumento levantado no levantamento de
literatura: há trabalho publicado sustentando que recuperação plana sobre
documento não estruturado limita raciocínio de múltiplos saltos quando os dados
são estruturados. Responder explicitamente, não por omissão.}

\subsubsection{Persistência e trilha de auditoria}

\pendente{Modelagem relacional, indexação de dados semiestruturados vindos da
API do PNCP, e versionamento com trilha que permita reproduzir uma pesquisa de
preços exatamente como ela era na data em que foi realizada. Este é o eixo de
maior interesse para a orientação e a lacuna mais defensável do trabalho.}

\pendente{Decisão entre persistência vetorial dedicada e a extensão do
PostgreSQL. A evidência é a Iniciação Científica em andamento no PET-EG, que
compara desempenho entre os dois caminhos. Critério acrescentado pela
orientação: verificar contra a natureza dos dados a tratar, não só o desempenho
medido. Se a decisão se apoiar em medição própria em vez de literatura, dizer
isso no texto --- não é evidência revisada por pares.}

\subsection{Desenvolvimento}

\pendente{Processo de desenvolvimento, com testes automatizados como critério
de conclusão de cada incremento.}

\subsection{Avaliação}

\pendente{Definir como a solução será avaliada e contra o quê. Sem uma linha de
base medida antes da intervenção, não há como afirmar redução de esforço ---
planejar a coleta dessa linha de base ainda em TCC I.}

\pendente{A etapa de verificação pelo órgão de controle ficou fora do escopo
obrigatório por decisão da orientação em 19/08/2026: o tempo de resposta do
órgão não é previsível e não cabe em planejamento. Reaproveitável no fim, se
houver tempo.}
